\documentclass[12pt,a4paper]{article}

\usepackage[spanish]{babel}
\usepackage[utf8]{inputenc}
\usepackage[T1]{fontenc}
\usepackage{geometry}
\usepackage{setspace}
\usepackage{graphicx}
\usepackage{tikz}
\usetikzlibrary{shapes.geometric, arrows.meta, babel} % Agregamos shapes y babel
\usepackage{enumitem}
\usepackage{titlesec}
\usetikzlibrary{shapes.geometric}

\geometry{margin=2.5cm}
\setstretch{1.5}

\title{Gestión de Procesos en las Organizaciones \\ \large Enfoque desde la Ingeniería de Sistemas}
\author{}
\date{}

\begin{document}

\maketitle

\section{Concepto de Proceso en una Organización}

\subsection{Definición}

Un proceso organizacional se define como un conjunto estructurado de actividades interrelacionadas que transforman entradas (inputs) en salidas (outputs), con el objetivo de generar valor para un cliente interno o externo.

Desde la perspectiva de la ingeniería de procesos, un proceso se caracteriza por:

\begin{itemize}
    \item Una secuencia lógica de actividades
    \item Recursos asignados (humanos, tecnológicos y materiales)
    \item Reglas y controles definidos
    \item Resultados medibles
\end{itemize}

En términos formales, un proceso es una unidad de trabajo que convierte recursos en productos o servicios mediante actividades coordinadas.

\subsection{Ejemplo Aplicado}

En el contexto de una empresa de comercio electrónico:

\begin{itemize}
    \item \textbf{Entrada:} Solicitud de compra del cliente
    \item \textbf{Actividades:} Validación del pedido, verificación de inventario, procesamiento de pago y despacho
    \item \textbf{Salida:} Producto entregado al cliente
\end{itemize}

Este proceso genera valor al satisfacer la necesidad del cliente de forma eficiente.

\subsection{Diagrama de Flujo del Proceso}

\begin{center}
\begin{tikzpicture}[node distance=2.5cm, auto]
    \node (start) [draw, ellipse] {Entrada};
    \node (val) [draw, rectangle, below of=start] {Validar pedido};
    \node (inv) [draw, rectangle, below of=val] {Verificar inventario};
    \node (pay) [draw, rectangle, below of=inv] {Procesar pago};
    \node (ship) [draw, rectangle, below of=pay] {Despacho};
    \node (end) [draw, ellipse, below of=ship] {Salida};

    \draw[->] (start) -- (val);
    \draw[->] (val) -- (inv);
    \draw[->] (inv) -- (pay);
    \draw[->] (pay) -- (ship);
    \draw[->] (ship) -- (end);
\end{tikzpicture}
\end{center}

\section{Tipos de Procesos en una Organización}

\subsection{Procesos Estratégicos}

Son aquellos que definen la dirección organizacional y están relacionados con la toma de decisiones a largo plazo. Establecen políticas, objetivos y estrategias.

\textbf{Ejemplo:} Planificación de innovación en una empresa tecnológica.

\subsection{Procesos Misionales}

También conocidos como procesos de negocio, están directamente vinculados con la razón de ser de la organización. Generan valor directo al cliente.

\textbf{Ejemplo:} Desarrollo de software en una empresa tecnológica.

\subsection{Procesos de Apoyo}

No generan valor directo, pero son esenciales para el funcionamiento de los procesos misionales.

\textbf{Ejemplo:} Gestión de recursos humanos.

\subsection{Relación entre Procesos}

\begin{center}
\begin{tikzpicture}[node distance=3cm]
    \node (est) [draw, rectangle] {Estratégicos};
    \node (mis) [draw, rectangle, below left of=est] {Misionales};
    \node (apo) [draw, rectangle, below right of=est] {Apoyo};

    \draw[->] (est) -- (mis);
    \draw[->] (est) -- (apo);
    \draw[<->] (mis) -- (apo);
\end{tikzpicture}
\end{center}

\section{Descomposición de Procesos}

\subsection{Definición General}

La descomposición de procesos consiste en dividir un proceso complejo en componentes más pequeños para facilitar su análisis, control y mejora.

\subsection{Elementos}

\subsubsection{Actividad}
Conjunto de tareas que representan una etapa del proceso.

\subsubsection{Tarea}
Unidad mínima de trabajo ejecutable.

\subsubsection{Procedimiento}
Descripción formal de cómo ejecutar actividades y tareas.

\subsubsection{Regla}
Condición que debe cumplirse para ejecutar el proceso.

\subsubsection{Norma}
Lineamiento formal que regula el proceso.

\subsection{Estructura Jerárquica}

\begin{center}
\begin{tikzpicture}[
    level 1/.style={sibling distance=5cm},
    level 2/.style={sibling distance=3cm},
    level 3/.style={sibling distance=2cm}
]
\node {Proceso}
    child {node {Actividades}
        child {node {Tareas}}
    }
    child {node {Procedimientos}}
    child {node {Reglas}}
    child {node {Normas}};
\end{tikzpicture}
\end{center}

\subsection{Ejemplo Integrado}

\begin{itemize}
    \item \textbf{Proceso:} Gestión de pedidos
    \item \textbf{Actividad:} Procesar pedido
    \item \textbf{Tarea:} Validar pago
    \item \textbf{Procedimiento:} Guía del sistema de pagos
    \item \textbf{Regla:} Pago aprobado
    \item \textbf{Norma:} Estándares de seguridad financiera
\end{itemize}

\section{Objetivos de un Proceso}

\subsection{Definición}

Los objetivos de un proceso son los resultados esperados, alineados con la estrategia organizacional. Deben ser:

\begin{itemize}
    \item Claros
    \item Medibles
    \item Alcanzables
    \item Relevantes
\end{itemize}

\subsection{Ejemplos}

\begin{itemize}
    \item Reducir tiempo de respuesta a menos de 24 horas
    \item Incrementar satisfacción del cliente al 90\%
    \item Disminuir errores operativos
\end{itemize}

\section{Límites de un Proceso}

\subsection{Definición}

Los límites determinan el inicio y fin de un proceso, delimitando su alcance.

\subsection{Componentes}

\begin{itemize}
    \item \textbf{Entrada:} Evento que inicia el proceso
    \item \textbf{Salida:} Resultado final
\end{itemize}

\subsection{Ejemplo}

\begin{itemize}
    \item Inicio: Cliente realiza compra
    \item Fin: Producto entregado
\end{itemize}

\section{Beneficios de la Gestión por Procesos}

\subsection{Eficiencia}
Optimización de recursos y reducción de tiempos.

\subsection{Control}
Monitoreo mediante indicadores.

\subsection{Mejora Continua}
Identificación de oportunidades de mejora.

\subsection{Estandarización}
Uniformidad en la ejecución de actividades.

\subsection{Reducción de Errores}
Minimización de fallos mediante reglas claras.

\section{Conclusión}

La gestión por procesos constituye un enfoque fundamental en la ingeniería de sistemas, ya que permite estructurar, analizar y optimizar el funcionamiento organizacional. Comprender la tipología, descomposición, objetivos y límites de los procesos facilita su modelado, especialmente en estándares como BPMN.

Este enfoque contribuye significativamente a la eficiencia operativa, la alineación estratégica y la mejora continua, aspectos esenciales en organizaciones modernas basadas en información y tecnología.

\end{document}
